% ===================================================================
\section{Conservation biology (EE07)}
% ===================================================================

\subsection{Why conserve: extinction rates and three ethics}
Contemporary extinction rates are \textbf{100--1000 times} the geological background, and projected to get worse.
Conservation biologists ``study the properties of small and/or declining populations and their environment, and try to
suggest measures to protect them in nature''; the field draws on population ecology, economics, public policy,
anthropology, law and more.

\begin{center}\small
\begin{tabularx}{\textwidth}{>{\raggedright\arraybackslash}p{3.6cm}X>{\raggedright\arraybackslash}p{2.9cm}}
\toprule
ethic (proponent) & premise & conclusion\\\midrule
romantic-transcendental (John Muir) & communion with nature brings people closer to God; morally superior to preserve
wilderness (pushed for Yosemite, Sequoia) & \textbf{preserve}\\
resource conservation (Gifford Pinchot; mentor Dietrich Brandis, first Inspector General of Forests in India) & ``the art
of producing from the forest whatever it can yield for the service of man''; wise use, greatest good for the greatest
number for the longest time (\emph{Bahujan Hitay, Bahujan Sukhai}) & \textbf{economic management} (sustainability)\\
evolutionary-ecological land ethic (Aldo Leopold) & ``a thing is right when it tends to preserve the integrity, stability
and beauty of the biotic community''; humans are plain members of the land community & \textbf{conserve} (use with
responsibility, intrinsic value)\\\bottomrule
\end{tabularx}
\end{center}
Different premises give different policies: no-go parks (Muir), sustained-yield forestry (Pinchot), managing for the
health of the whole community (Leopold).

\subsection{Who is vulnerable: Rabinowitz's seven rarities}
Cross three binary properties: \textbf{geographic range} (extensive/restricted), \textbf{habitat tolerance}
(broad/narrow), \textbf{local population size} (large/small). $2^3=8$ cells: one is \emph{common} (extensive, broad, large:
house sparrow, dandelion), seven are forms of rarity, ending with the rarest (restricted, narrow, small: mountain
gorilla).
\begin{itemize}
\item \textbf{Passenger pigeon}: extensive range, \emph{narrow} habitat (mast forests), \emph{large} local populations:
``abundant but demanding''.
\item \textbf{Tiger}: extensive range, broad tolerance, \emph{small} local populations: ``everywhere but sparse''.
\end{itemize}
Each road to rarity calls for a different remedy (protect corridors and prey; protect the key habitat across the range;
guard the last populations). ``Rare'' describes a state; vulnerability also depends on threats.

\subsection{The passenger pigeon: a complicated story}
The last bird died in \textbf{1914}, but \textbf{1878} was the last year the catch was commercially viable. Hunting
pressure fell, yet the population never rebounded, though other pigeons routinely do. Slide 12's six candidate causes:
climate/catastrophe (no records); \textbf{hunting} (part of the story, not the whole); disease (no records);
competition with domestic species (unlikely); \textbf{habitat alteration coupled with social facilitation in food
finding} (``we do not know''); \textbf{genetic issues} (Murray et al.\ 2017 found low genomic diversity shaped by
selection; how it contributed: ``we have no idea'').

\begin{chain}[: why a flock species can be doomed while still common (Bucher, background)]
\begin{enumerate}
\item The pigeon was a \textbf{nomadic specialist} on mast (beechnuts, acorns), which appears in large but unpredictable
patches; finding them depended on \textbf{social facilitation} (following flock-mates).
\item Deforestation, pigs eating mast, and hunting made patches smaller and rarer, and fewer birds meant poorer
food finding.
\item Poorer food finding lowered per-capita success, which shrank the flocks further: a \textbf{positive feedback}.
\item So per-capita growth \emph{fell} with falling density: an \textbf{Allee effect}. The MVP was crossed while the bird
was still relatively common; after that, removing the hunting pressure could not bring it back.
\end{enumerate}
\end{chain}

\subsection{Why small populations are at risk}
Three stochasticities (slide 15):
\begin{itemize}
\item \textbf{demographic}: random variation in births, deaths, sex ratio;
\item \textbf{genetic}: drift, loss of heterozygosity, inbreeding depression;
\item \textbf{environmental} stochasticity and natural catastrophes.
\end{itemize}
\textbf{Minimum viable population (MVP)}: the smallest size that ensures, at an acceptable level of risk, that the
population persists for a specified time (conventionally 99\,\% over 1000 years, Shaffer 1981).

\begin{chain}[: why size matters for each stochasticity]
\begin{itemize}
\item \emph{Demographic}: relative noise scales as $1/\sqrt N$ (32\,\% at $N=10$, 3\,\% at $N=1000$). Even with $b>d$, each
founder's line dies out with probability $d/b$, and all must fail: $P_{\text{ext}}=(d/b)^{N_0}>0$. With $b/d=1.2$: one
founder fails 83\,\% of the time, 20 founders 2.6\,\%. Mean persistence time grows \emph{exponentially} with $K$.
\item \emph{Genetic}: $H_t=H_0(1-1/2N_e)^t$; the half-life is $\approx1.39N_e$ generations. Lost variation means
inbreeding depression and less capacity to adapt.
\item \emph{Environmental}: hits every individual together, so it does \emph{not} average out with $N$; large
populations are much safer, not safe.
\end{itemize}
\end{chain}

\textbf{Allee effect}: per-capita growth falls as density falls (mates hard to find, group defence and food-finding fail).
Strong Allee model: $\dot N=rN(N/A-1)(1-N/K)$, with equilibria 0 (stable), $A$ (unstable threshold), $K$ (stable). A
constant harvest can create a second, higher threshold: with $r=0.5$, $A=20$, $K=100$ and quota 20, the population
sits at 85.9 but a bad year that drops it below 52.1 sends it to extinction; at quota $\ge26.3$ the equilibria collide
(a fold, as in lakes).

\begin{keyidea}[: the extinction vortex]
The three threats compound rather than add: smaller $N$ $\Rightarrow$ faster loss of heterozygosity $\Rightarrow$
inbreeding lowers survival and fecundity $\Rightarrow$ lower $b/d$ at smaller $N$ $\Rightarrow$ steeply higher extinction
risk; in social species the Allee effect adds a further drop. Once a population is below the size where these feedbacks
dominate, removing the original pressure does not save it. That is the pigeon after 1878.
\end{keyidea}

\textbf{Four things must converge} (slide 16): science, common sense, economics and politics.

\subsection{Project Tiger}
\begin{itemize}
\item \textbf{Baseline}: 40--45 thousand tigers in early-20th-century India; \textbf{1{,}827 in 1972} (hunting, habitat
loss, prey loss, poaching, competition with villagers and livestock).
\item \textbf{Launch 1 April 1973}. Aims: a viable tiger population for scientific, economic, aesthetic, cultural and
ecological values; preserve such areas as national heritage.
\item \textbf{Site criteria}: little disturbance; high numbers of prey and predator; no drilling, mining or timber;
spread across states. \textbf{First nine reserves} ($\sim$13{,}017\,km$^2$): Bandipur, Corbett, Kanha, Manas, Melghat,
Palamau, Ranthambhore, Similipal, Sundarbans.
\item \textbf{Activities}: core and buffer areas, corridors, halting cattle and forest harvesting, patrolling.
\item \textbf{Result}: $>$4{,}000 in 11 years (questioned, but an increase was real).
\item \textbf{Reversals}: lack of political will against farmer pressure, buffer encroachment, prey loss, poaching, lack of
funds. More reserves did not mean more tigers.
\item \textbf{Reform}: Wild Life (Protection) Amendment Act 2006 created the \textbf{NTCA} and the \textbf{Wildlife Crime
Control Bureau}; camera-trap surveys; 1{,}411 (2008), 1{,}706 (2011); Sariska lost all its tigers (poaching); 2012 Supreme
Court ban on tourists in core areas; \textbf{2018: 2{,}967 (2{,}603--3{,}346)} from 2{,}461 unique photographs.
\item India holds $\sim$\textbf{70\,\%} of the world's wild tigers. CITES (1973) regulates trade. Human--animal conflict:
elephants and railways; lemon-tree barriers in Assam.
\end{itemize}

\begin{chain}[: why counts need a detection probability]
A raw count $C$ misses animals that were present but not seen. If each animal is detected with probability $p$, the
estimate is $\hat N=C/\hat p$. Indices (pugmarks, scats) ignore $p$, and the pugmark ``census'' could not even identify
individuals reliably (prints vary with substrate and pace); it was abandoned. Camera traps identify tigers by
\textbf{stripes}; a short (30--60 day) survey keeps the population \emph{closed}; capture--recapture models estimate $p$
from how often known individuals are re-photographed. Lincoln--Petersen: $\hat N=Mn/m$ (marked $M$, second sample
$n$, recaptures $m$), from $m/n=M/N$. Fewer recaptures, wider confidence interval.
\end{chain}

\begin{worked}[: 2025 mid-sem Q3]
46 rabbits marked; next week 37 caught, 27 marked. $\hat N=46\times37/27=63$ (Chapman: $47\cdot38/28-1=62.8$).
\end{worked}

\subsection{Reserves (mandatory reading: Hunter and Gibbs Ch.\ 11)}
Reserves are the \textbf{coarse-filter} strategy: protect representative ecosystems. \textbf{Selection}: diversity
hotspots, gap filling, logistics. \textbf{Diamond's six principles} (reserves as islands): (1) large holds more species
than small; (2) one large better than several small of equal total area; (3) small reserves close together;
(4) clustered rather than linear; (5) corridors; (6) as circular as possible.

\begin{chain}[: why large and round]
Area grows as the square of linear size, edge only linearly; so larger reserves have relatively less edge, and for a
given area a circle has the least perimeter. Edge is where poachers, livestock, fire, dogs, disease and invaders enter,
and where animals leave and meet conflict. Large reserves also hold more habitat types, wide-ranging species, larger
(safer) populations, and cannot all be hit by one catastrophe. They should exceed the \textbf{minimum dynamic area}
(room for a mosaic of disturbance and succession patches).
\end{chain}

\textbf{SLOSS} (single large or several small): with \emph{nested} species lists (larger reserves contain the smaller ones'
species) one large reserve wins (224 vs 199 species for the same cost); with \emph{un-nested} lists several small ones win
(709 species for less money). The answer depends on how species are distributed.

\textbf{Connectivity}: boundaries are permeable; buffers of semi-natural land help. Four movements to maintain: daily
(home range), annual migration, dispersal of young and seeds, long-term range shifts. Corridors and \textbf{zoning} (core,
buffer, partial use). \textbf{Management} ``is first and foremost people management'': visitors, legal uses, local
goodwill; let natural disturbances (fire, flood) act; control invasives and overabundant natives.

\begin{statements}
\begin{enumerate}[label=\arabic*.]
\item ``A population whose birth rate exceeds its death rate cannot go extinct.''
\F $P_{\text{ext}}=(d/b)^{N_0}>0$; small populations often do.
\item ``Since hunting ended around 1878, the passenger pigeon's extinction in 1914 must have had a cause unrelated to hunting.''
\F Hunting (with habitat loss) pushed it below a threshold; after that, Allee-type feedbacks drove the decline without
needing a new cause.
\item ``Environmental stochasticity becomes negligible in large populations, like demographic stochasticity.''
\F It hits all individuals together and does not average out.
\item ``A species with an extensive range and broad tolerance is safe from extinction.''
\F The tiger is in that cell but has small local populations: demographic and genetic risk, Allee effects (2025 end-sem Q3).
\item ``A five-pointed star reserve of the same area is as good as a circle.''
\F More edge, narrow arms that are almost all edge, smaller core: more poaching, conflict, invasion (2025 end-sem Q7).
\item ``If the buffer zone degrades, tigers in an untouched core are unaffected.''
\F Buffers hold prey, dispersal routes for subadults, and keep people and cattle away from the core (2025 end-sem Q9).
\item ``Several small reserves can protect more species than one large one of equal cost.''
\T When species lists are un-nested; SLOSS has no universal answer.
\item ``A rise in the number of tiger reserves shows that tiger numbers rose.''
\F Reserves measure effort, not outcome; reversals happened while reserves increased.
\item ``A constant harvest below the maximum growth rate is always safe.''
\F With an Allee effect, harvesting creates an unstable threshold far above $A$; a bad year can tip the stock.
\end{enumerate}
\end{statements}

\begin{longans}[: ``Why are small populations at risk even when healthy on average?'' / ``Project Tiger as a case study'']
\emph{Small populations}: three stochasticities with the $1/\sqrt N$ argument and $(d/b)^{N_0}$; drift and inbreeding;
environmental shocks; Allee effect; how they compound (vortex); MVP definition; passenger pigeon as the example.
\emph{Project Tiger}: baseline and causes; aims, criteria, design (core, buffer, corridors); early success and why it
was questioned (pugmarks); reversals; 2006 reform; camera-trap capture--recapture and $\hat N=C/\hat p$; numbers to 2018;
lessons: science, common sense, economics, politics, and local people (conflict).
\end{longans}
